\section{Variance bound for 3 colorings of expanders}

% some setup
% remove things that have been discussed before

\providecommand{\Vbad}{V_{\mathrm{bad}}}

Let~\(G\) be an undirected graph with vertex set~\(V\).

Let \(\bm x\bm y\) denote a uniformly random edge of \(G\).
Since the graph is undirected, the random variables \(\bm x\) and \(\bm y\) have the same distribution without loss of generality.
For a vertex subset \(S\subseteq V\), we denote by \(\mu(S)=\Pr\{\bm x\in S\}\) its probability under this distribution.
(For regular graphs, this probability is the fraction of vertices contained in \(S\).)

Let \(A\) denote the random walk (transition) matrix for \(G\).
In other words, \(A\) is the adjacency matrix of \(G\) with every column normalized to sum to \(1\).
This matrix acts as a linear operator on real-valued functions \(f\from V\to \R\),
\[
A f (x) = \E \Brac{f(\bm y) \Mid \bm x = x}\mper
\]
This operator is self-adjoint with respect to the inner product for functions \(f,g\from V\to \R\),
\[
\iprod{f,g} \seteq \E f(\bm x) g(\bm x)\mper
\]
We denote the Euclidean norm with respect to this inner product by \(\norm{f}\seteq \iprod{f,f}^{1/2}=\Paren{\E f^{2}}^{1/2}\).

The linear operator \(A\) has largest eigenvalue \(1\) with all-ones \(\Ind\) as eigenfunction.
Let \(\lambda\) denote the second largest eigenvalue of \(A\).
If \(G\) has a spectral gap, i.e., \(\lambda<1\), all other eigenfunctions \(\xi\) are orthogonal to \(\Ind\) with respect to the above inner product and thus satisfies \(\E \xi=0\).
Every function \(f\from V\to \R\) has projection \((\E f)\cdot \Ind\) to the all-ones eigenspace and thus satisfies the inequality,
\begin{equation}
  \iprod{f,Af}\le (\E f)^{2} + \lambda \cdot\Paren{\E f^{2}- \Paren{\E f}^{2}}\mper
\end{equation}

Suppose~\(\chi\from V\to [k]\) is a (not necessarily proper) \(k\)-coloring of~\(G\).
Let~\(\Vbad\) be a vertex cover of monochromatic edges under~\(\chi\).
For \(a\in [k]\), let \(\chi_{a}\from V\to \set{0,1}\) denote the indicator function of color class \(a\).
We let \(\psi_{a}\from V\to \R_{\ge 0}\) denote the image of \(\chi_{a}\) under \(A\) so that  \(\psi_{a}=A\chi_{a}\) and,
\begin{equation}
  \psi_{a}(x) = \Pr\Set{\chi(\bm y)= a \Mid \bm x = x}\mper
\end{equation}

\begin{lemma}
  For all distinct \(a,b\in [k]\), the variance of \(\psi_{b}\) conditioned on \(\chi_{a}\) satisfies
  \begin{equation}
    \E\Brac{ (\psi_{b}(\bm x)-\alpha_{a,b})^{2}\Mid \chi(\bm x)=a}
    \le \lambda \tfrac{\E \chi_{b}}{\E \chi_{a}} + \lambda^{-1} E_{a,b}
    \mcom
  \end{equation}
  where \(\alpha_{a,b}=\E\Brac{\psi_{b}(\bm x)\Mid \chi(\bm x)=a}\) and \(E_{a,b}\lesssim \mu\Paren{\Vbad \cap \chi^{-1}(a)} + \lambda^{2} \cdot \mu\Paren{\Vbad\cap \chi^{-1}(b)}\).
  \david{more natural to bound the error in terms of the fraction of monochromatic edges as opposed to the size of the vertex cover.}
\end{lemma}

\begin{proof}
  For scalars \(\alpha,\beta>0\) to be determined later, we consider the function \(f= \chi_{a} \cdot (\psi_{b}-\alpha) + \beta\cdot \chi_{b}\).
  The expectation of this function satisfies \(\E f =  \E \chi_{a} \cdot (\psi_{b}-\alpha) + \beta \E\chi_{b}\).
  Since the functions \(\chi_{a}\cdot \psi_{b}\) and \(\chi_{b}\) are orthogonal, we can compute the norm of \(f\) as,
  \begin{equation}
    \E f^{2}= \E \chi_{a}\cdot(\psi_{b}-\alpha)^{2} + \beta^{2}\E\chi_{b}
    \mper
  \end{equation}
  Since \(\Vbad\) covers all monochromatic edges, we can bound the quadratic form of \(A\) evaluated at \(f\),
  \begin{align}
    \iprod{f,Af}
    &=  \iprod{\chi_{a} \cdot (\psi_{b}-\alpha),A (\chi_{a} \cdot (\psi_{b}-\alpha))} + \beta^{2}\iprod{\chi_{b},A\chi_{b}} + 2\beta\cdot \iprod{\chi_{a}(\psi_{b}-\alpha), A\chi_{b}}\\
    &\ge -2 \mu\Paren{\Vbad \cap \chi^{-1}(a)} - 2 \beta^{2} \cdot \mu\Paren{\Vbad\cap \chi^{-1}(b)} + 2\beta \cdot\E \chi_{a} (\psi_{b}-\alpha)^{2}
    \mper
  \end{align}
  Here, we also use that \(A\) is an expectation operator so that \(\norm{A g}_{\infty}\le \norm{g}_{\infty}\) for every function \(g\from V\to \R\) and that \(\E \chi_{a}(\psi_{b}-\alpha)=0\).
  
  By choosing \(\alpha=\E\chi_{a}\psi_{b} \cdot \Paren{\E \chi_{a}}^{-1}\), we cancel the first term in our expression for \(\E f\) so that this expectation simplifies to \(\E f=\beta \E \chi_{b}\).
  We abbreviate the above bound on the contribution of monochromatic edges by \(E_{a,b}\seteq 2 \mu\Paren{\Vbad \cap \chi^{-1}(a)} + 2 \beta^{2} \cdot \mu\Paren{\Vbad\cap \chi^{-1}(b)} \).
  In this way, we obtain the inequality,
  \begin{align}
    2\beta \E \chi_{a} \cdot(\psi_{b}-\alpha)^{2}
    & \le E_{a,b} + \iprod{f,Af} \\
    & \le E_{a,b} +  (\E f)^{2} + \lambda \E f^{2} \\
    & =   E_{a,b} + \beta^{2} (\E \chi_{b})^{2} + \lambda \E \chi (\psi_{b}-\alpha)^{2} + \lambda \beta^{2} \E \chi_{b}
  \end{align}
  Since the variance term \(\E \chi_{a} \cdot (\psi_{b}-\alpha)^{2}\) appears on both sides of the inequality with different coefficients, we can obtain an upper bound on this term by choosing \(\beta=\lambda\) and solving the resulting inequality for the variance term,
  \begin{equation}
    \E \chi_{a} \cdot(\psi_{b}-\alpha)^{2}
    \le \lambda^{-1} E_{a,b} + \lambda (\E \chi_{b})^{2} + \lambda^{2} \E\chi_{b}
    \le\lambda^{-1} E_{a,b} + 2\lambda \E\chi_{b}
    \mcom
  \end{equation}
  as desired.
\end{proof}


%%% Local Variables:
%%% mode: LaTeX
%%% TeX-master: "../main"
%%% End:
